% =====================================================================
% Report content to be pasted INSIDE \begin{document} ... \end{document}
%
% Required packages in your Overleaf preamble:
%   \usepackage{graphicx}
%   \usepackage{amsmath, amssymb}
%   \usepackage{float}        % for [H] figure placement
%   \usepackage{booktabs}     % nice tables (\toprule, \midrule, \bottomrule)
%   \usepackage{listings}     % code snippets
%   \usepackage{xcolor}
%   \usepackage{hyperref}     % cross-references
%   \lstset{basicstyle=\ttfamily\footnotesize, breaklines=true,
%           keywordstyle=\color{blue}, commentstyle=\color{gray!80!black},
%           numbers=left, numberstyle=\tiny, stepnumber=1, frame=single,
%           language=Matlab}
%
% Place the folders 'images_part1' and 'images_part2' next to your main.tex
% =====================================================================

\title{Lateral Vehicle Dynamics and Lane-Keeping Control \\ \large Project 2 -- Driver Assistance System Design}
\author{Dennis [Surname] \\ Politecnico di Torino}
\date{\today}
\maketitle

%======================================================================
\section{Introduction}
%======================================================================
The lateral dynamics of a road vehicle govern its directional behaviour
during cornering and overtaking, and ultimately determine the
effectiveness of any lane-keeping or path-tracking driver-assistance
function. The linear single-track (``bicycle'') model captures the
two main degrees of freedom of interest -- sideslip angle $\beta$ and
yaw rate $\dot\psi$ -- with sufficient accuracy for stability and
controller-design tasks, while the dynamic single-track (DST) and the
nonlinear Pacejka variant (DSTP) are used as more realistic plants for
closed-loop simulations.

This project is organised in two parts:
\begin{itemize}
    \item \textbf{Part 1 -- Lateral Dynamics.} Computation of the
          derivatives of stability, root-locus analysis of the
          linearised state-space system as a function of speed, time
          response to a step steer, and steady-state vehicle gains
          (path curvature, lateral acceleration, sideslip).
    \item \textbf{Part 2 -- Lane-Keeping Control.} Implementation of
          DST/DSTP models in Simulink, design and comparison of
          PDF/PIDF controllers (with two error structures) at
          $65$\,km/h and $100$\,km/h, implementation of a Stanley
          kinematic controller and final benchmark of all the
          strategies.
\end{itemize}

The same vehicle parameters are used throughout Part 1, summarised in
Table~\ref{tab:params}. The cornering stiffness values $C_1$ and
$C_2$ are not constant input data: they are extracted from the
\texttt{TireData.xlsx} file (Pacejka block) by linear fitting of the
$F_y(\alpha)$ curve near zero slip, evaluated at the static axle loads
$F_{z,1}\!\approx\!6$\,kN (front) and $F_{z,2}\!\approx\!4$\,kN (rear)
that follow from the longitudinal mass distribution.

\begin{table}[H]
    \centering
    \begin{tabular}{lcl}
        \toprule
        \textbf{Symbol} & \textbf{Value} & \textbf{Description} \\
        \midrule
        $a$    & $0.923$\,m       & Front axle -- c.o.g. distance \\
        $b$    & $1.376$\,m       & Rear axle -- c.o.g. distance \\
        $m$    & $1020$\,kg       & Vehicle mass \\
        $J_z$  & $1225$\,kg\,m$^2$& Yaw inertia \\
        $t$    & $1.365$\,m       & Track width \\
        $C_1$  & from tyre file   & Front-axle cornering stiffness \\
        $C_2$  & from tyre file   & Rear-axle cornering stiffness \\
        \bottomrule
    \end{tabular}
    \caption{Vehicle parameters used in Part 1.}
    \label{tab:params}
\end{table}

%======================================================================
\section{Lateral Dynamics -- Theoretical Background}
\label{sec:theory}
%======================================================================
Under the assumptions of no aerodynamic forces, nil self-aligning
torque at the wheel contact and a linear tyre characteristic, the
equations of high-speed cornering reduce to
\begin{align}
    mV(\dot\beta + \dot\psi) + m\dot V\beta &=
        Y_\beta\,\beta + Y_{\dot\psi}\,\dot\psi + Y_\delta\,\delta,
        \label{eq:lat_y}\\
    J_z\,\ddot\psi &=
        N_\beta\,\beta + N_{\dot\psi}\,\dot\psi + N_\delta\,\delta,
        \label{eq:lat_n}
\end{align}
where the \emph{derivatives of stability} are
\begin{equation}
\begin{aligned}
    Y_\beta      &= -(C_1 + C_2),                     &
    N_\beta      &= -(a C_1 - b C_2), \\
    Y_{\dot\psi} &= -\tfrac{a C_1 - b C_2}{V},        &
    N_{\dot\psi} &= -\tfrac{a^2 C_1 + b^2 C_2}{V}, \\
    Y_\delta     &= C_1,                              &
    N_\delta     &= a\,C_1 .
\end{aligned}
\label{eq:derivatives}
\end{equation}
With the state vector $\mathbf{x} = [\beta,\,\dot\psi]^T$ and input
$\delta$, neglecting external forces and braking torques, the
state-space form reads
\begin{equation}
    \dot{\mathbf{x}} = A\,\mathbf{x} + B\,\delta, \quad
    A =
    \begin{bmatrix}
        \dfrac{Y_\beta}{mV} & \dfrac{Y_{\dot\psi}}{mV} - 1 \\[6pt]
        \dfrac{N_\beta}{J_z} & \dfrac{N_{\dot\psi}}{J_z}
    \end{bmatrix},
    \quad
    B =
    \begin{bmatrix}
        \dfrac{Y_\delta}{mV} \\[6pt] \dfrac{N_\delta}{J_z}
    \end{bmatrix}.
    \label{eq:ss}
\end{equation}

\paragraph{Steady-state gains.}
By zeroing the derivatives in~\eqref{eq:lat_y}--\eqref{eq:lat_n} and
solving for the four classical performance ratios one gets
\begin{align}
    \frac{1}{R\delta}    &= \frac{Y_\delta N_\beta - N_\delta Y_\beta}
        {V\bigl[N_\beta(mV - Y_{\dot\psi}) + N_{\dot\psi} Y_\beta\bigr]},
        \label{eq:gain_curv}\\[2pt]
    \frac{V^2}{R\delta}  &= V\,\frac{Y_\delta N_\beta - N_\delta Y_\beta}
        {N_\beta(mV - Y_{\dot\psi}) + N_{\dot\psi} Y_\beta},
        \label{eq:gain_ay}\\[2pt]
    \frac{\beta}{\delta} &= \frac{-N_\delta(mV - Y_{\dot\psi}) - N_{\dot\psi} Y_\delta}
        {N_\beta(mV - Y_{\dot\psi}) + N_{\dot\psi} Y_\beta},
        \label{eq:gain_beta}
\end{align}
which respectively express the path curvature, the lateral
acceleration and the sideslip angle per unit steering input.

%======================================================================
\section{Part 1 -- Lateral Dynamics}
%======================================================================

\subsection{Cornering stiffness extraction}
The MATLAB script \texttt{part\_1.m} loads the Pacejka block of
\texttt{TireData.xlsx} (columns Q:V, rows 3:203), which contains the
lateral force $F_y$ versus slip angle $\alpha$ at five vertical loads
$F_z \in \{2,4,6,8,10\}$\,kN. The cornering stiffness per tyre is
identified as the slope $-\mathrm{d}F_y/\mathrm{d}\alpha$ near
$\alpha = 0$, evaluated at the static axle loads
\begin{equation}
    F_{z,1} = \tfrac{m g b}{a+b} \approx 5990\,\text{N}, \qquad
    F_{z,2} = \tfrac{m g a}{a+b} \approx 4015\,\text{N},
\end{equation}
which justifies the use of the columns corresponding to $6$\,kN and
$4$\,kN. Doubling the per-tyre value (two tyres per axle) gives the
bicycle-model equivalents $C_1$ and $C_2$:
\begin{lstlisting}
idx = abs(alpha_deg) < 2;            % points close to alpha = 0
p6  = polyfit(alpha(idx), Fy(idx,3), 1);  Cp_front = -p6(1);  % Fz = 6 kN
p4  = polyfit(alpha(idx), Fy(idx,2), 1);  Cp_rear  = -p4(1);  % Fz = 4 kN
C1  = 2*Cp_front;                    % front axle
C2  = 2*Cp_rear;                     % rear  axle
\end{lstlisting}

\subsection{Stability derivatives and state-space at 80 km/h}
With $V_0 = 80/3.6$\,m/s and the identified stiffness, the local
function \texttt{bicycleSS} returns the numerical derivatives of
stability and the matrices $A$, $B$ in~\eqref{eq:ss}:
\begin{lstlisting}
function [A,B,der] = bicycleSS(V,C1,C2,a,b,m,Jz)
    der.Yb = -(C1 + C2);
    der.Yr = -(a*C1 - b*C2)/V;
    der.Nb = -(a*C1 - b*C2);
    der.Nr = -(a^2*C1 + b^2*C2)/V;
    der.Yd = C1;
    der.Nd = a*C1;
    A = [ der.Yb/(m*V),  der.Yr/(m*V)-1 ;
          der.Nb/Jz,     der.Nr/Jz      ];
    B = [ der.Yd/(m*V) ;
          der.Nd/Jz    ];
end
\end{lstlisting}
At $80$\,km/h the eigenvalues of $A$ are a stable complex-conjugate
pair, consistent with the lightly damped yaw/sideslip oscillation
typical of a passenger car at moderate-to-high speed. The negative
real part guarantees asymptotic stability of the free response.

\subsection{Root loci versus speed}
Sweeping $V$ in the interval $[2,60]$\,m/s and computing the
eigenvalues of $A(V)$ traces the locus of the two closed-loop poles
of the lateral dynamics. The result is shown in
Figure~\ref{fig:rootloci}.

\begin{figure}[H]
    \centering
    \includegraphics[width=0.85\textwidth]{images_part1/Root loci vs speed.png}
    \caption{Root loci of the $[\beta,\,\dot\psi]$ state-space system
    as a function of vehicle speed.}
    \label{fig:rootloci}
\end{figure}

At very low speed the two poles lie on the negative real axis and are
far apart ($\sim\!-160$ and $\sim\!-8$), denoting an overdamped
heading response dominated by tyre friction. As the speed increases,
the two real poles approach each other, collide on the real axis at
the breakaway point ($V\!\approx\!12$--$14$\,m/s, i.e. $\sim 45$
km/h) and then become a complex-conjugate pair. The imaginary part
grows with $V$ -- reaching $\pm 7$\,rad/s at the upper end of the
sweep -- while the real part remains negative and stabilises around
$-3\,\div\,-5$, showing a progressive reduction of damping ratio with
speed. The locus stays entirely in the left half plane: the open-loop
vehicle is stable over the full speed range explored.

\subsection{Step steer at 80 km/h}
A small step steer command $\delta_0 = 0.03$\,rad ($\approx 1.7^\circ$
at the wheel) is applied to the linear state-space system at
$80$\,km/h. The time response of $\beta$ and $\dot\psi$ is plotted in
Figure~\ref{fig:steepsteer}.

\begin{figure}[H]
    \centering
    \includegraphics[width=0.9\textwidth]{images_part1/steep_steer.png}
    \caption{Step steer $\delta = 0.030$\,rad at $80$\,km/h:
    sideslip angle $\beta$ (top) and yaw rate $\dot\psi$ (bottom).}
    \label{fig:steepsteer}
\end{figure}

The yaw-rate response is well damped: $\dot\psi$ overshoots the
steady-state value of $\sim 0.20$\,rad/s by about $5\%$ at
$t\!\approx\!0.25$\,s and then settles within half a second.
The sideslip $\beta$ exhibits the classical non-minimum-phase-like
behaviour of an understeering vehicle at high speed: it rises
positively to a peak of $\approx 4\!\times\!10^{-3}$\,rad and then
crosses zero, settling at a small \emph{negative} value of
$\approx -2\!\times\!10^{-3}$\,rad. Physically, the front tyres
develop lateral force first, the chassis rotates outwards, and only
later the rear axle catches up giving a negative steady-state
sideslip -- a clear signature of mild understeer.

\subsection{Steady-state gains}
The four gain expressions~\eqref{eq:gain_curv}--\eqref{eq:gain_beta}
are evaluated over the same speed range as the root locus.
Figure~\ref{fig:gains} reports curvature, lateral acceleration and
sideslip angle gains versus speed in km/h, with a vertical dashed
line marking the design speed of $80$\,km/h.

\begin{figure}[H]
    \centering
    \includegraphics[width=0.95\textwidth]{images_part1/Figure_3.png}
    \caption{Steady-state gains versus vehicle speed: path curvature
    $1/(R\delta)$ (top), lateral acceleration $V^2/(R\delta)$ (middle)
    and sideslip angle $\beta/\delta$ (bottom).}
    \label{fig:gains}
\end{figure}

The three plots confirm the qualitative classification of the
vehicle:
\begin{itemize}
    \item the \textbf{curvature gain} decreases monotonically with
          speed -- from $\approx 0.43$\,m$^{-1}$ at low speed down to
          $\approx 0.10$\,m$^{-1}$ at $200$\,km/h. The vehicle needs
          progressively more steering input to negotiate a given
          curve as it speeds up, a hallmark of an \emph{understeering}
          characteristic;
    \item the \textbf{lateral acceleration gain} grows almost
          linearly with $V$ up to $\sim 80$\,km/h and then bends
          downwards as the saturation $V^2/R\delta$ approaches its
          asymptote. At $80$\,km/h the gain is $\approx 150$
          (m/s$^2$)/rad, i.e. $0.26$\,g per degree of steering;
    \item the \textbf{sideslip angle gain} crosses zero at
          $\approx 80$\,km/h, becoming negative at higher speed.
          This crossover is the \emph{characteristic speed} predicted
          by the understeer theory: above it, the steady-state
          sideslip points outwards from the curve centre, a
          condition that requires the driver to add steering input
          to maintain the trajectory.
\end{itemize}
The design speed of $80$\,km/h sits exactly on the sideslip
crossover, which makes it a natural reference point for the
controller-design phase in Part 2.

%======================================================================
\section{Part 2 -- Lane-Keeping Control}
%======================================================================
The second part of the project shifts the focus from open-loop
analysis to closed-loop lane keeping. A reference trajectory
$Y_r(X) = 10\sin(0.01\,X)$ -- a sine wave with $10$\,m amplitude and
$\approx 628$\,m wavelength -- is sampled at $20\,000$ points; the
target heading $\psi_r$ is obtained by numerical differentiation. The
MATLAB driver \texttt{part2.m} asks the user for the cruise speed
(allowed values: $65$ or $100$\,km/h) and chains three Simulink
models, one per exercise. The simulation horizon is $T_\mathrm{fin} =
150$\,s.

%----------------------------------------------------------------------
\subsection{Exercise 1 -- DST and DSTP open-loop models}
%----------------------------------------------------------------------

\subsubsection*{Simulink implementation}
Two distinct vehicle plants are implemented in
\texttt{project\_2\_part2\_ex1.slx}: a Dynamic Single-Track (DST)
model with a \emph{linear} tyre characteristic and a Dynamic
Single-Track with \emph{Pacejka} tyre (DSTP), see
Figure~\ref{fig:sim_ex1_inner}. In both models the steer input goes
through a first-order actuator $\dfrac{1}{0.1 s + 1}$, the differential
equations of the DST/DSTP block are implemented in a MATLAB Function,
and the six states $[X, Y, \psi, v_x, v_y, \omega_\psi]^T$ are
integrated by a single $1/s$ block. The outer wrapper
(Figure~\ref{fig:sim_ex1_outer}) routes the same sinusoidal input
$\delta_f(t) = 0.5\sin(0.2\,t)$\,rad to both plants and exports the
relevant signals to the workspace.

\begin{figure}[H]
    \centering
    \includegraphics[width=0.95\textwidth]{images_part2/ex1/simulink_DST_DSTP_models.png}
    \caption{Internal structure of the DST (linear) and DSTP
    (Pacejka) Simulink blocks used in Exercise 1.}
    \label{fig:sim_ex1_inner}
\end{figure}

\begin{figure}[H]
    \centering
    \includegraphics[width=0.95\textwidth]{images_part2/ex1/simulink_1.png}
    \caption{Top-level wrapper of Exercise 1: the same steering input
    drives both models in parallel and the signals are exported to
    MATLAB.}
    \label{fig:sim_ex1_outer}
\end{figure}

\subsubsection*{Simulation results}
With initial conditions
$(X_0,Y_0)=(0,0)$\,m, $\psi_0 = 0^\circ$,
$v_x(0) = 80$\,km/h, $v_y(0) = \omega_\psi(0) = 0$, no longitudinal
acceleration ($a_x = 0$) and the sinusoidal steering input above,
the two plants produce the trajectories shown in
Figure~\ref{fig:ex1_results}.

\begin{figure}[H]
    \centering
    \includegraphics[width=\textwidth]{images_part2/ex1/DST_DSTP_plots.png}
    \caption{Exercise 1: trajectory $(X,Y)$ and longitudinal speed
    $v_x$ for the linear DST model (top row) and the nonlinear DSTP
    model (bottom row).}
    \label{fig:ex1_results}
\end{figure}

The linear model produces a quasi-periodic trajectory bounded inside
a $\sim\!100\times 80$\,m box: the lateral force scales linearly
with slip angle, so each steering cycle generates the same yaw
moment and the vehicle traces a series of overlapping loops. The
longitudinal speed drops sharply from $22$\,m/s to a steady
$\approx 11$\,m/s and oscillates only mildly, because the linear tyre
keeps reabsorbing energy through the $v_y\omega_\psi$ coupling in
$\dot v_x$.

The Pacejka model, conversely, drifts to the bottom right of the
plane reaching $X\!\approx\!250$\,m and $Y\!\approx\!-180$\,m. The
saturated lateral force at large slip angles limits the yaw moment
the vehicle can develop, so the trajectory is no longer symmetric and
the longitudinal speed settles around $10.5$\,m/s after a slower
transient. This is the expected qualitative difference: the linear
tyre overestimates the cornering capability of the vehicle and
underestimates its longitudinal energy loss, while the Pacejka tyre
delivers a realistic, asymmetric, partly-saturating behaviour.

%----------------------------------------------------------------------
\subsection{Exercise 2 -- PDF and PIDF lane-keeping controllers}
%----------------------------------------------------------------------

\subsubsection*{Error definitions and controller structure}
Two error signals are extracted by the \texttt{errors} MATLAB
Function from the reference and the actual front-axle pose:
\begin{itemize}
    \item the \textbf{cross-track error} $e_{ct}$, i.e. the lateral
          distance of the front axle from the reference path,
          projected onto the local normal direction;
    \item the \textbf{heading error} $e_h = \psi_a - \psi_r$, i.e.
          the angular mismatch between vehicle and tangent.
\end{itemize}
For each cruise speed two controller architectures are tested:
\begin{enumerate}
    \item input = $e_{ct}$ only (the heading is ignored by the
          controller and recovers passively through the closed loop);
    \item input = $e_{ct} + e_h$ (a simple sum of the two errors as
          a coarse stand-in for a feed-forward heading term).
\end{enumerate}
The plant is the DSTP nonlinear model, fed through a first-order
$1/(0.1 s + 1)$ steering actuator and instrumented with the
\texttt{translatePose} block, which translates the centre-of-gravity
pose to the front-axle pose used by the controller
(Figure~\ref{fig:sim_ex2_plant}).

\begin{figure}[H]
    \centering
    \includegraphics[width=\textwidth]{images_part2/ex2/simulink_DSTP_and_translatePose.png}
    \caption{Vehicle block used in Exercise 2: DSTP model with
    steering actuator and \texttt{translatePose} for front-axle
    output.}
    \label{fig:sim_ex2_plant}
\end{figure}

Two complete loops are run in parallel inside
\texttt{project\_2\_part2\_ex2.slx} -- one with input $e_{ct}$
(Figure~\ref{fig:sim_ex2_ect}) and one with input $e_{ct} + e_h$
(Figure~\ref{fig:sim_ex2_ect_eh}). Both feed a PDF or a PIDF block
(selected by the workspace flags \texttt{controller1}/\texttt{controller2}),
tuned through the Simulink PID Tuner around the linearisation point
$(X,Y,\psi)=(0,0,\psi_r(1))$, $v_x = v_{x,l}$,
$v_y=\omega_\psi=0$.

\begin{figure}[H]
    \centering
    \includegraphics[width=\textwidth]{images_part2/ex2/simulink_e_ct_PIDF_PDF.png}
    \caption{Closed loop with controller input $e_{ct}$ only.}
    \label{fig:sim_ex2_ect}
\end{figure}

\begin{figure}[H]
    \centering
    \includegraphics[width=\textwidth]{images_part2/ex2/simulink_e_ct_e_h_PIDF_PDF.png}
    \caption{Closed loop with controller input $e_{ct} + e_h$.}
    \label{fig:sim_ex2_ect_eh}
\end{figure}

\subsubsection*{Results at 65 km/h}
At the lower speed both controllers succeed in keeping the vehicle
on the sinusoidal reference. Among the six combinations available
(two errors $\times$ two controllers, plus the two steering-angle
traces), the most representative are reported below; the remaining
ones (PDF with $e_{ct}$ alone, and the PDF steering signals at
$65$\,km/h) follow the same trend with negligible additional
information and are summarised numerically in Section~\ref{sec:bench}.

\begin{figure}[H]
    \centering
    \includegraphics[width=0.95\textwidth]{images_part2/ex2/65_km_h/e_ct_PIDF.png}
    \caption{PIDF with input $e_{ct}$, $v_x = 65$\,km/h.}
    \label{fig:ex2_65_pidf_ect}
\end{figure}

\begin{figure}[H]
    \centering
    \includegraphics[width=0.95\textwidth]{images_part2/ex2/65_km_h/e_ct_e_h_PIDF.png}
    \caption{PIDF with input $e_{ct}+e_h$, $v_x = 65$\,km/h.}
    \label{fig:ex2_65_pidf_ect_eh}
\end{figure}

\begin{figure}[H]
    \centering
    \includegraphics[width=0.95\textwidth]{images_part2/ex2/65_km_h/steering_angle_PIDF.png}
    \caption{Steering angle commanded by the PIDF controller at
    $65$\,km/h, with input $e_{ct}$ (top) and $e_{ct}+e_h$ (bottom).}
    \label{fig:ex2_65_steer_pidf}
\end{figure}

At $65$\,km/h the PIDF with $e_{ct}$
(Figure~\ref{fig:ex2_65_pidf_ect}) tracks the reference within a
$\pm 0.035$\,m envelope, with the heading error confined below
$0.1^\circ$. Switching to $e_{ct}+e_h$
(Figure~\ref{fig:ex2_65_pidf_ect_eh}) does not improve the steady
performance significantly: the cross-track error grows to
$\pm 0.05$\,m because the controller now devotes part of its action
to the heading channel as well. The PDF counterparts behave similarly
with errors approximately three times larger ($\sim 0.11$\,m and
$\sim 0.16$\,m respectively), due to the absence of the integral
action that compensates for the slow drift introduced by the
discretisation noise in $\dot Y_r$. The steering angle
(Figure~\ref{fig:ex2_65_steer_pidf}) oscillates between $\pm 0.25^\circ$,
i.e. comfortably below the physical actuator limit of $\delta_m =
35^\circ$.

\subsubsection*{Results at 100 km/h -- the divergence of PIDF$(e_{ct}+e_h)$}
At $100$\,km/h the picture changes drastically. The same controllers
tuned around the linearisation point start to interact with the
unmodelled fast lateral dynamics, and -- because the linearisation
was performed at a different speed -- the closed-loop poles move
towards the right of the complex plane.

\begin{figure}[H]
    \centering
    \includegraphics[width=0.95\textwidth]{images_part2/ex2/100_km_h/e_ct_PDF.png}
    \caption{PDF with input $e_{ct}$, $v_x = 100$\,km/h: large but
    transient excursion, then good tracking.}
    \label{fig:ex2_100_pdf_ect}
\end{figure}

\begin{figure}[H]
    \centering
    \includegraphics[width=0.95\textwidth]{images_part2/ex2/100_km_h/e_ct_PIDF.png}
    \caption{PIDF with input $e_{ct}$, $v_x = 100$\,km/h: very large
    cross-track excursion ($\sim 60$\,m) but eventual recovery.}
    \label{fig:ex2_100_pidf_ect}
\end{figure}

\begin{figure}[H]
    \centering
    \includegraphics[width=0.95\textwidth]{images_part2/ex2/100_km_h/e_ct_e_h_PIDF.png}
    \caption{PIDF with input $e_{ct}+e_h$, $v_x = 100$\,km/h:
    \emph{unstable} closed loop. The cross-track error grows above
    $900$\,m and the heading oscillates with $\pm 180^\circ$ swings.}
    \label{fig:ex2_100_pidf_ect_eh}
\end{figure}

\begin{figure}[H]
    \centering
    \includegraphics[width=0.95\textwidth]{images_part2/ex2/100_km_h/steering_angle_PIDF.png}
    \caption{Steering angle of the PIDF controller at $100$\,km/h:
    saturation at $\sim 600^\circ$ on the top trace and a runaway up
    to $3.5\!\times\!10^4$\,deg on the bottom trace
    ($e_{ct}+e_h$ case).}
    \label{fig:ex2_100_steer_pidf}
\end{figure}

The PDF with $e_{ct}$ alone (Figure~\ref{fig:ex2_100_pdf_ect})
survives the transient: after a $\sim 10$\,m cross-track deviation
around $t = 12$\,s, induced by the higher lateral acceleration that
the controller was not tuned for, the loop recovers within $\sim 30$\,s
and stays on path. The PIDF with $e_{ct}$ alone
(Figure~\ref{fig:ex2_100_pidf_ect}) is much worse: the integral wind-up
amplifies the initial mismatch and the cross-track error grows to
$\sim\!60$\,m before slowly returning towards zero.

The catastrophic case is the PIDF with $e_{ct}+e_h$
(Figure~\ref{fig:ex2_100_pidf_ect_eh}): the controller never
recovers. After $\sim 30$\,s the cross-track error breaks loose and
grows monotonically beyond $900$\,m, while the heading error
oscillates with $\pm 180^\circ$ swings and the commanded steering
angle (Figure~\ref{fig:ex2_100_steer_pidf}, bottom trace) reaches
$3.5\!\times\!10^{4}$\,deg, far beyond any physical actuator. The
root cause is twofold: (i) the sum $e_{ct}+e_h$ couples two
quantities with very different sensitivities to speed -- the
contribution of $e_h$ doubles when $V$ doubles -- so what was a
benign feed-forward term at $65$\,km/h becomes a destabilising loop
gain at $100$\,km/h; (ii) the integral channel of the PIDF
accumulates the resulting bias, and the wind-up brings the closed
loop to instability.

\paragraph{Take-away of Exercise 2.}
At low speed all four controllers (PDF/PIDF with one of the two error
structures) provide acceptable tracking, and the PIDF with $e_{ct}$
alone gives the smallest steady error. At high speed the same tuning
is no longer robust: only the PDF with $e_{ct}$ stays usable, while
the PIDF with $e_{ct}+e_h$ diverges. This motivates the introduction
of a controller whose structure is intrinsically speed-aware -- the
Stanley kinematic controller of Exercise 3.

%----------------------------------------------------------------------
\subsection{Exercise 3 -- Stanley kinematic controller}
%----------------------------------------------------------------------

\subsubsection*{Control law}
The Stanley controller computes the steering command as
\begin{equation}
    \delta_S = \mathrm{sat}\!\left(
        e_h + \arctan\!\frac{k\,e_{ct}}{v_b + v_x}\,
    \right),\qquad |\delta_S| \le \delta_m,
    \label{eq:stanley}
\end{equation}
where $v_b = 1$\,m/s is a small speed bias that prevents the
denominator from vanishing at low speed and $\delta_m = 35^\circ$ is
the physical actuator limit. The first term aligns the heading with
the local path tangent; the second term forces a cross-track
correction whose strength is inversely proportional to the actual
speed -- which is precisely the speed-awareness that the PIDF
lacked.

\begin{figure}[H]
    \centering
    \includegraphics[width=\textwidth]{images_part2/ex3/simulink_stanely.png}
    \caption{Simulink implementation of the Stanley closed loop
    (\texttt{project\_2\_part2\_ex3.slx}).}
    \label{fig:sim_ex3}
\end{figure}

The Simulink loop (Figure~\ref{fig:sim_ex3}) reuses the
\texttt{ref\_gen} + \texttt{errors} chain from Exercise 2 and replaces
the PDF/PIDF block by a MATLAB Function implementing
equation~\eqref{eq:stanley}.

\subsubsection*{Comparison at 65 and 100 km/h}

\begin{figure}[H]
    \centering
    \includegraphics[width=0.95\textwidth]{images_part2/ex3/comparison_PIDF_Stanley_65.png}
    \caption{PIDF vs Stanley at $65$\,km/h.}
    \label{fig:ex3_65}
\end{figure}

\begin{figure}[H]
    \centering
    \includegraphics[width=0.95\textwidth]{images_part2/ex3/comparison_PIDF_Stanley_100.png}
    \caption{PIDF vs Stanley at $100$\,km/h.}
    \label{fig:ex3_100}
\end{figure}

At $65$\,km/h (Figure~\ref{fig:ex3_65}) the PIDF wins on cross-track
accuracy ($\pm 0.035$\,m against $\pm 0.16$\,m), but the heading
error of Stanley is appreciably smoother (a clean sinusoidal trace
of the same amplitude as the PIDF $\pm 0.10^\circ$ noisy signal).
The reason is that Stanley reads the heading error directly and
produces a noise-free $\delta$, while the PIDF amplifies the
high-frequency content of the numerical $\psi_r$.

At $100$\,km/h (Figure~\ref{fig:ex3_100}) the contrast becomes
decisive: the PIDF develops the same $\sim\!60$\,m transient already
discussed, whereas Stanley sticks to the reference with a cross-track
error in the centimetre range and a heading error well below
$0.1^\circ$. The arctangent saturation effectively limits the
authority of the cross-track channel and prevents the runaway
observed with PIDF$(e_{ct}+e_h)$.

\paragraph{Take-away of Exercise 3.}
Stanley trades a slightly larger steady-state cross-track at low
speed for a much more graceful behaviour at high speed and zero
sensitivity to controller-tuning speed mismatches. For a real
lane-keeping function operating across a wide speed range, this is
the safer trade-off.

%======================================================================
\section{Controllers Benchmark}
\label{sec:bench}
%======================================================================
The numerical performance of every (controller, error, speed)
combination is summarised in
Table~\ref{tab:bench}. The figures are read from the time histories
of $e_{ct}(t)$ and $e_h(t)$ within the simulation horizon
$T_\mathrm{fin}=150$\,s; ``stable'' means the error stays bounded and
the steady-state tracking is preserved, ``transient'' means a finite
excursion is followed by recovery, ``unstable'' means the error grows
unbounded until the steering saturates.

\begin{table}[H]
    \centering
    \footnotesize
    \begin{tabular}{lllrrl}
        \toprule
        \textbf{Controller} & \textbf{Input} & \textbf{Speed} &
        \textbf{$|e_{ct}|_\text{max}$} & \textbf{$|e_h|_\text{max}$} &
        \textbf{Behaviour} \\
        \midrule
        PDF     & $e_{ct}$        & $65$\,km/h  & $\sim 0.11$ m   & $\sim 0.07^\circ$  & stable \\
        PIDF    & $e_{ct}$        & $65$\,km/h  & $\sim 0.035$ m  & $\sim 0.08^\circ$  & stable, best at 65 \\
        PDF     & $e_{ct}+e_h$    & $65$\,km/h  & $\sim 0.16$ m   & $\sim 0.10^\circ$  & stable \\
        PIDF    & $e_{ct}+e_h$    & $65$\,km/h  & $\sim 0.05$ m   & $\sim 0.10^\circ$  & stable \\
        Stanley & $e_{ct},\,e_h$  & $65$\,km/h  & $\sim 0.16$ m   & $\sim 0.10^\circ$  & stable, smooth $\delta$ \\
        \midrule
        PDF     & $e_{ct}$        & $100$\,km/h & $\sim 10$ m     & $\sim 80^\circ$    & transient, recovers \\
        PIDF    & $e_{ct}$        & $100$\,km/h & $\sim 60$ m     & $\sim 180^\circ$   & transient, slow recovery \\
        PDF     & $e_{ct}+e_h$    & $100$\,km/h & $\sim 0.4$ m    & $\sim 3^\circ$     & stable, growing oscillation \\
        PIDF    & $e_{ct}+e_h$    & $100$\,km/h & $> 900$ m       & $\pm 180^\circ$    & \textbf{unstable} \\
        Stanley & $e_{ct},\,e_h$  & $100$\,km/h & $< 0.05$ m      & $< 0.1^\circ$      & stable, best at 100 \\
        \bottomrule
    \end{tabular}
    \caption{Benchmark of the lane-keeping controllers for the
    sinusoidal reference $Y_r(X)=10\sin(0.01\,X)$.}
    \label{tab:bench}
\end{table}

Three observations stand out:
\begin{itemize}
    \item At $65$\,km/h the \textbf{PIDF with $e_{ct}$} is the most
          accurate option, with sub-centimetre tracking, but Stanley
          comes very close on the heading channel and with a much
          cleaner steering command.
    \item Increasing the speed to $100$\,km/h destroys every linear
          controller tuned at the lower speed. The most striking
          collapse is PIDF$(e_{ct}+e_h)$, whose steering command runs
          away to $3.5\!\times\!10^4$\,deg. A controller designed at
          $65$\,km/h \emph{cannot} be safely re-used at $100$\,km/h
          without re-tuning or gain scheduling.
    \item Stanley is the only controller that retains good tracking
          across the entire speed range, thanks to its explicit
          dependence on $v_x$ in the cross-track term.
\end{itemize}

%======================================================================
\section{Conclusions}
%======================================================================
The project covered the full pipeline of lateral vehicle modelling
and control, starting from the linear bicycle model and ending with
a non-linear closed-loop lane-keeping simulation.
\begin{itemize}
    \item \textbf{Part 1} confirmed that the bicycle model with
          tyre-data-derived cornering stiffness reproduces the
          essential features of an understeering passenger car:
          stable pole locus over the full speed range, well-damped
          step-steer response, decreasing curvature gain and a
          sideslip-gain sign change at the characteristic speed of
          $\sim 80$\,km/h.
    \item \textbf{Exercise 1 of Part 2} highlighted the qualitative
          difference between the DST (linear tyre) and DSTP (Pacejka
          tyre) models: under the same aggressive sinusoidal steer,
          the linear model produces a quasi-periodic trajectory while
          the Pacejka one drifts due to tyre saturation. The DSTP is
          therefore the appropriate plant for the controller-design
          stage.
    \item \textbf{Exercise 2} showed that, at the design speed
          $65$\,km/h, both PDF and PIDF achieve sub-metric tracking on
          a sinusoidal lane. Increasing the speed to $100$\,km/h
          without re-tuning the gains exposes a clear lack of
          robustness: the PIDF with $e_{ct}+e_h$ becomes unstable,
          the PIDF with $e_{ct}$ alone develops a $60$\,m transient,
          while the PDF with $e_{ct}$ recovers from a $\sim\!10$\,m
          excursion.
    \item \textbf{Exercise 3} demonstrated that the Stanley
          kinematic controller, by embedding the speed $v_x$ in the
          denominator of its cross-track term, achieves a uniform
          high-quality tracking across $65$ and $100$\,km/h. It pays
          a small price in steady-state cross-track at low speed but
          remains the only architecture viable at high speed without
          re-tuning.
\end{itemize}

The benchmark in Section~\ref{sec:bench} synthesises the message of
the project: \emph{linear controllers tuned around a single operating
point are intrinsically speed-sensitive, while a kinematic controller
that takes the longitudinal speed explicitly into account preserves
its performance over a wider envelope}. A natural extension of this
work is a gain-scheduled PIDF (tuned at multiple speeds) or a
model-predictive lane-keeping controller, both of which would
combine the steady-state accuracy of PIDF at low speed with the
robustness of Stanley at high speed.
